\scanpage{31}
section $\lambda$ de $\mathbb G_m$, donc \deleted{\unsure{$\alpha$|$*$}}

\[
\alpha=\lambda+\lambda^{-1},
\]

et on est ramené à une question de calcul
standart.

\deleted{Pour \unsure{analyser}}

\uline{Cor. 1} Supposons les car. résiduelles de $S\ne3$. Pour qu'
une section $u$ de $G$ satisfasse $u^3=\mathrm{id}$ (resp. $u^6=\mathrm{id}$) et que
l'ordre de $u_s$ ($s\in S$) soit tjrs égal à 3 (resp 6) il f et s.
qu'on ait $\alpha=-1$ (resp. $\alpha=+1$).

\deleted{\unsure{L'une et l'autre hyp. signifie que|L'une et l'autre hyp. implique que}}

\uline{Cor 2} Supposons les car. rés. de $S\ne2$. Pour que l'on
ait $u^4=\mathrm{id}$ et que les $u_s$ $(s\in S)$ soient d'ordre 4 exactement,
il f. et s. que $\alpha=0$

\uline{Cor. 3} Supposons les car. résiduelles $\ne5$. Pour que
$u^5=0$ et les $u_s$ d'ordre 5, il f et s. qu'on ait

\[
\alpha^2+\alpha-1=0
\]

\uline{NB} L'ordre \inserted{$\nu$} d'un élément $\tilde u$ de $\mathrm{Sl}(2,k)$ ($k$ un corps)
-- donc d'un $\widetilde G(k)$ pour une forme \inserted{$\widetilde G$} de $\mathrm{Sl}(2)$ \unsure{sur|sur $k$}
\deleted{\unsure{ou}} peut être d'une des formes suivantes \deleted{\unsure{ou $p$}}
pour une car. donnée

a) \inserted{$\nu=$} $+\infty$\\ 
b) $\nu$ entier $\geq1$ \uline{premier à} $p$ :\\ 
c) $\nu=2p$ ($p$ premier \uuline{impair})\\ 
\deleted{\unsure{multiple de $p$}}\\ 
d) $\nu=p$
\margin{si $p=\mathrm{car.}\,k$ est \unsure{finie|impair|fini} $\Big\}$}

Si on exclut les cas $\pm\mathrm{id}$ (élémt du centre, i.e.
d'un élément qui donne $1$ dans $\widetilde G/\mu_2=G$)
[qui donnent lieu à $\alpha=\pm2$, mais d'autres cas y
donnent lieu aussi] on a vu que la valeur
de $\alpha$ détermine (loc. fppf) la classe de conj. de
$u$. On trouve en fait l'analyse suivante :

% archive sheet number: 30
